\documentclass[a4paper,10pt]{amsart}

\pdfinfoomitdate=1
\pdftrailerid{}
\pdfsuppressptexinfo=15

\usepackage[T1]{fontenc}
\usepackage{lmodern}
\usepackage[margin=27mm]{geometry}
\usepackage{amsmath,amssymb,mathtools,booktabs,array,microtype,xcolor}
\usepackage[shortlabels]{enumitem}
\usepackage[numbers,sort&compress]{natbib}
\usepackage[colorlinks=true,linkcolor=blue!55!black,citecolor=blue!55!black,urlcolor=blue!55!black]{hyperref}
\hypersetup{pdftitle={},pdfauthor={},pdfsubject={},pdfkeywords={}}

\newtheorem{theorem}{Theorem}
\newtheorem{proposition}[theorem]{Proposition}
\newtheorem{lemma}[theorem]{Lemma}
\newtheorem{corollary}[theorem]{Corollary}
\theoremstyle{remark}
\newtheorem{remark}[theorem]{Remark}

\newcommand{\F}{\mathbb F}
\newcommand{\cP}{\mathcal P}
\newcommand{\Fix}{\operatorname{Fix}}
\newcommand{\imof}{\operatorname{im}}

\title[Steiner triangle collapse]{A Four-Stratum Functional Graph and Exact Fibres for a Binary-Projective Steiner Triangle Map}
\author{Anonymous}
\date{}

\begin{document}

\begin{abstract}
Let $\cP=\F_2^r\setminus\{0\}$ and endow it with the standard Steiner
quasigroup operation $x\star x=x$ and $x\star y=x+y$ for $x\ne y$.
We determine the complete functional graph of the ordered-triple map
\[
 (a,b,c)\longmapsto(b\star c,c\star a,a\star b).
\]
Diagonal triples and projective-line blocks are fixed, triples with exactly
two equal entries form strict $3$-cycles, and every remaining triple has depth
one.  If $N=2^r-1$, the fixed, strict-$3$-cycle, and depth-one counts are
$N^2$, $N(N-1)$, and $N(N-1)(N-3)$, respectively.  More sharply, every
target fibre has size $1$, $N-2$, or $0$ according as the target is diagonal
or two-equal, a distinct block, or a nonblock; block sources admit an explicit
one-parameter description.  We also identify the symmetry orbits and weak
components, compute all fibre moments, and recover $r$ from the unlabelled
fibre multiset.  The Steiner system, its quasigroup operation,
and characteristic-two projective model are established inputs and receive
no contribution credit.  An exact audit checks the claimed graph and inverse
formulas through rank six.
\end{abstract}

\maketitle

\section{Owned inputs and the finite map}\label{sec:setup}

Fix $r\ge2$, put $V=\F_2^r$, and let
\[
 \cP=V\setminus\{0\},\qquad N=|\cP|=2^r-1.
\]
The points and lines of $\operatorname{PG}(r-1,2)$ form the classical
projective Steiner triple system: the line through distinct $x,y\in\cP$ is
$\{x,y,x+y\}$~\cite{JungnickelTonchev2019,FalconeFigulaGalici2025}.
The associated Steiner quasigroup is
\begin{equation}\label{eq:star}
 x\star y=
 \begin{cases}
 x,&x=y,\\
 x+y,&x\ne y.
 \end{cases}
\end{equation}
The correspondence between Steiner triple systems and Steiner quasigroups is
standard~\cite{DrapalKozlikGriggs2012}.  Quasigroup and Steiner local rules
also occur in cellular automata~\cite{Moore1997}.  All of these facts are
background here.

We study the finite self-map
\begin{equation}\label{eq:S}
 S:\cP^3\longrightarrow\cP^3,
 \qquad S(a,b,c)=(b\star c,c\star a,a\star b).
\end{equation}
The recent ``Steiner products and dynamics'' of
Kettinger--Peterson~\cite{KettingerPeterson2025} concern an
orientation-dependent anticommutative bilinear product on a real vector space
and its fixed-vector iterates.  That is not the coordinate map~\eqref{eq:S}.
This title-level neighbour is included to make the source boundary explicit,
not to convert a bounded search non-hit into a priority claim.

\begin{table}[ht]
\centering
\small
\caption{Source subtraction.  The middle column is established input; the
last column is the residual analyzed here.}
\label{tab:subtraction}
\begin{tabular}{@{}>{\raggedright\arraybackslash}p{.22\textwidth}
                    >{\raggedright\arraybackslash}p{.34\textwidth}
                    >{\raggedright\arraybackslash}p{.34\textwidth}@{}}
\toprule
Source area & Zero-credit input & Retained special-map output \\
\midrule
Steiner systems and quasigroups
& operation~\eqref{eq:star}, projective point-line blocks, standard identities
& no carrier or quasigroup claim \\
Quasigroup cellular automata
& algebraic local rules and their predictability
& finite ordered-triple graph, not a lattice CA \\
Oriented Steiner products
& a real bilinear product and fixed-vector iteration
& coordinatewise map~\eqref{eq:S}, four strata, and target fibres \\
\bottomrule
\end{tabular}
\end{table}

Call a triple \emph{diagonal} if all entries agree, \emph{two-equal} if
exactly two agree, a \emph{block} if its entries are distinct and
$a+b+c=0$, and a \emph{nonblock} otherwise.  These four classes partition
$\cP^3$.

\begin{theorem}[complete graph and target fibres]\label{thm:main}
For the map $S$ in~\eqref{eq:S} the following statements hold.
\begin{enumerate}[(i)]
\item Every diagonal and every block is fixed.  Every two-equal triple has
exact period three.  Every nonblock has depth one and maps to a fixed block.
\item There are $N^2$ fixed points, $N(N-1)$ strict $3$-cycles, and
$N(N-1)(N-3)$ depth-one states.  Hence
\begin{equation}\label{eq:periodic}
 |\imof S|=|\operatorname{Per}S|=4N^2-3N.
\end{equation}
The maximum tail is zero for $r=2$ and one for $r\ge3$.
\item For every target $q\in\cP^3$,
\begin{equation}\label{eq:fibres}
 |S^{-1}(q)|=
 \begin{cases}
 1,&q\text{ is diagonal or two-equal},\\
 N-2,&q\text{ is a block},\\
 0,&q\text{ is a nonblock}.
 \end{cases}
\end{equation}
For a block $q=(x,y,z)$ its sources are exactly
\begin{equation}\label{eq:param}
 (a,b,c)=(t,t+z,t+y),\qquad t\in V\setminus\{0,y,z\}.
\end{equation}
The same fibre formula holds for $S^k$ for every $k\ge1$.
\item The fixed-point counts and Artin--Mazur zeta function are
\begin{equation}\label{eq:fixk}
 |\Fix(S^k)|=
 \begin{cases}N^2,&3\nmid k,\\4N^2-3N,&3\mid k,
 \end{cases}
 \qquad
 \zeta_S(q)=\frac1{(1-q)^{N^2}(1-q^3)^{N(N-1)}}.
\end{equation}
\end{enumerate}
\end{theorem}

\begin{proposition}[symmetry, components, and fibre moments]\label{prop:structure}
The map is equivariant for the natural action of
$\operatorname{GL}(r,2)\times\mathfrak S_3$.  For $r\ge3$, the four classes
above are exactly its four orbits on $\cP^3$; for $r=2$ the nonblock orbit is
empty.  The functional graph has
\begin{equation}\label{eq:components}
 2N^2-N
\end{equation}
weak components: $N$ isolated diagonal fixed points, $N(N-1)$ isolated
$3$-cycles, and $N(N-1)$ rooted stars whose centers are block fixed points
and whose $N-3$ leaves are nonblocks.  In particular, for every integer
$j\ge1$,
\begin{equation}\label{eq:moments}
 \sum_{q\in\cP^3}|S^{-1}(q)|^j
 =3N^2-2N+N(N-1)(N-2)^j.
\end{equation}
For $N>3$ the maximum fibre is $N-2$ and is attained exactly at blocks; the
proportion of Garden-of-Eden states is
$(N-1)(N-3)/N^2$, which tends to one with $r$.
\end{proposition}

\begin{corollary}[fibre enumerator and rank recovery]\label{cor:fibrepoly}
The codomain fibre enumerator is
\begin{equation}\label{eq:fibrepoly}
 \Phi_r(u):=\sum_{q\in\cP^3}u^{|S^{-1}(q)|}
 =N(N-1)(N-3)+(3N^2-2N)u+N(N-1)u^{N-2}.
\end{equation}
Terms with equal exponents are understood to combine; in particular, when
$r=2$ every target has fibre one.  For every $r\ge2$, the unlabelled fibre
multiset alone recovers the projective rank:
\begin{equation}\label{eq:recover}
 r=\log_2\bigl(\max_q|S^{-1}(q)|+3\bigr).
\end{equation}
For a uniformly chosen codomain target, the probability of being in the image
is $4/N-3/N^2$, whereas the Garden-of-Eden probability tends to one.
\end{corollary}

The retained result is the conjunction of the four-stratum graph and the
target-resolved inverse law.  Neither the projective design nor the operation
\eqref{eq:star} is claimed.

For orientation, Table~\ref{tab:ranks} displays consequences at the first
three ranks.  Rank two has no transient shell; rank three is the first case in
which most states have no predecessor.

\begin{table}[ht]
\centering
\small
\caption{Exact consequences of Theorem~\ref{thm:main}.  The ``$3$-cycles''
column counts cycles rather than states.}
\label{tab:ranks}
\begin{tabular}{@{}rrrrrr@{}}
\toprule
$r$ & $N$ & states & fixed & $3$-cycles & depth one \\
\midrule
2 & 3  & 27   & 9   & 6   & 0 \\
3 & 7  & 343  & 49  & 42  & 168 \\
4 & 15 & 3375 & 225 & 210 & 2520 \\
\bottomrule
\end{tabular}
\end{table}

\section{The four strata}\label{sec:strata}

\begin{proof}[Proof of Theorem~\ref{thm:main}(i)]
Idempotence makes every $(x,x,x)$ fixed.  If $x\ne y$, direct use of
\eqref{eq:star} gives, for one of the three possible equality positions,
\begin{equation}\label{eq:threecycle}
 (x,x,y)\mapsto(x+y,x+y,x)\mapsto(y,y,x+y)\mapsto(x,x,y).
\end{equation}
The three displayed states are distinct, so the period is exactly three.
The other equality positions behave by the same cyclic coordinate symmetry.

For pairwise distinct $a,b,c$, no idempotent branch is used, and
\begin{equation}\label{eq:linearbranch}
 S(a,b,c)=(b+c,c+a,a+b).
\end{equation}
These three outputs are nonzero and pairwise distinct: a zero output or an
equality would contradict distinctness of the inputs.  Their sum is zero.
Thus every distinct triple maps to a block.  If $a+b+c=0$, then
$b+c=a$, $c+a=b$, and $a+b=c$, so the block is fixed.  A nonblock cannot be
fixed and therefore has depth exactly one.
\end{proof}

\begin{proof}[Proof of Theorem~\ref{thm:main}(ii)]
There are $N$ diagonals.  An ordered block is determined by an ordered pair
of distinct first entries, because the third is their nonzero sum.  Hence
there are $N(N-1)$ blocks and $N^2$ fixed points in total.

There are $3N(N-1)$ two-equal states.  Equation~\eqref{eq:threecycle}
partitions them into $N(N-1)$ strict cycles.  The number left after removing
the diagonal, two-equal, and block strata from $N^3$ is
\[
 N^3-N-3N(N-1)-N(N-1)=N(N-1)(N-3).
\]
These are precisely the depth-one nonblocks.  The image contains the entire
periodic locus.  It contains nothing else by the classification, proving
\eqref{eq:periodic}.  The depth-one count vanishes exactly when $N=3$, which
is the boundary rank $r=2$.
\end{proof}

The boundary is worth emphasizing.  On the three-point projective line every
distinct ordered triple is already a block, so small-rank inspection alone
would hide the collapsing shell.  The first Garden-of-Eden states occur in
the Fano carrier $r=3$.

\section{A direct inverse parametrization}\label{sec:inverse}

\begin{proof}[Proof of Theorem~\ref{thm:main}(iii)]
The equality calculation in~\eqref{eq:threecycle} shows directly that each
diagonal or two-equal target has its unique predecessor in the same stratum.
Equation~\eqref{eq:linearbranch} shows that a distinct source can reach only
a block; hence a nonblock target has no source.

It remains to invert a block $(x,y,z)$.  Any distinct source must solve
\begin{equation}\label{eq:system}
 b+c=x,\qquad c+a=y,\qquad a+b=z.
\end{equation}
The block condition $x+y+z=0$ is exactly the consistency relation.  Set
$a=t$.  Then~\eqref{eq:system} gives the family~\eqref{eq:param}.
Its entries are all nonzero exactly when $t\notin\{0,z,y\}$.  They are then
automatically distinct: for example $t=t+z$ would force $z=0$, and the other
two equalities similarly force $y=0$ or $x=y+z=0$.  Conversely every
allowed $t$ yields a valid source.  Since $|V|=N+1$, there are $N-2$ sources.

Among them, $t=x$ recovers the fixed block itself; the remaining $N-3$ are
nonblocks and have no predecessors.  Thus passing from $S$ to any positive
iterate neither gains nor loses a source for a block.  The same is immediate
on the two permutation strata, proving the last assertion.
\end{proof}

The inverse proof is independent of division by the orbit census: it solves
the target equations and exposes the three forbidden parameter values.

\begin{proof}[Proof of Proposition~\ref{prop:structure}]
An invertible linear map preserves equality and sends $x+y$ to the sum of
the images.  A coordinate permutation merely permutes the three opposite
pair products.  These observations prove equivariance.  Diagonals are one
linear orbit.  A two-equal triple is an ordered independent pair together
with an equality position, blocks are ordered independent pairs completed by
their sum, and nonblocks are ordered independent triples.  Transitivity of
$\operatorname{GL}(r,2)$ on ordered independent lists, together with the
coordinate action, gives the orbit claim.

Theorem~\ref{thm:main} shows that a diagonal has no attached transient, a
$3$-cycle has no attached transient, and each block has exactly $N-3$
nonblock sources in addition to itself.  This proves~\eqref{eq:components}
and the component description.  There are $N+3N(N-1)=3N^2-2N$ targets of
fibre one and $N(N-1)$ block targets of fibre $N-2$; all others have fibre
zero.  This gives~\eqref{eq:moments}.  The remaining assertions follow from
the nonblock count in Theorem~\ref{thm:main}(ii).
\end{proof}

\begin{proof}[Proof of Corollary~\ref{cor:fibrepoly}]
Insert the three target counts used in the proof of
Proposition~\ref{prop:structure} into the definition of $\Phi_r$.  When
Theorem~\ref{thm:main}(iii) gives maximum fibre $N-2$; hence
$N=\max|S^{-1}(q)|+2$ and $2^r=N+1$, proving~\eqref{eq:recover}.  Finally
divide~\eqref{eq:periodic} by $N^3$.  Its complement is the nonblock ratio
already stated in Proposition~\ref{prop:structure}.
\end{proof}

\begin{proof}[Proof of Theorem~\ref{thm:main}(iv)]
Only fixed points and strict $3$-cycles are recurrent.  Therefore a recurrent
state is fixed by $S^k$ precisely as stated in~\eqref{eq:fixk}.  Substitution
in the defining exponential identity
$\zeta_S(q)=\exp(\sum_{k\ge1}|\Fix(S^k)|q^k/k)$, or equivalently multiplication
over primitive cycles, gives the displayed product.
\end{proof}

\section{Exact audit and limitations}\label{sec:audit}

The accompanying standard-library verifier enumerates every state through
$r=6$.  For each rank it compares the literal update with the four-way
classification, orbit depth and period, every target indegree, the explicit
parameter set~\eqref{eq:param}, all counts, and fixed points of several
iterates.  This is finite falsification pressure and regression control, not
an all-rank proof or an ownership certificate.

The result is specific to the binary projective Steiner quasigroup and to
ordered triples with the simultaneous pairwise update~\eqref{eq:S}.  It does
not classify arbitrary Steiner quasigroups or triple systems, assert stability
under perturbing the rule, or furnish a general quasigroup-dynamics theorem.
The source search was bounded and terminology dependent; no novelty, priority,
or external-release claim is made.

\subsection*{Data availability}
No empirical data were used.  The exact verifier and frozen transcript are
included with the manuscript.

\subsection*{Ethics, contributions, conflicts, and funding}
No human or animal subjects are involved.  Authorship is anonymous for
internal review.  The authors declare no conflict of interest and no external
funding for this work.

\bibliographystyle{plainnat}
\bibliography{references}

\end{document}
